\documentclass[12pt]{article}
\usepackage[utf8]{inputenc}
\usepackage[T1]{fontenc}
\usepackage{hyperref}
\title{Minjungbal Language Resources for K-12 Curriculum Development}
\author{Compiled by Claude Sonnet 5.5 (claude-sonnet-5-5) for Andr\'as Kornai}
\date{}
\begin{document}
\maketitle

\section{Introduction}
Minjungbal (ISO 639-3 xjb; Glottocode twee1234) is a Pama-Nyungan variety of the Tweed Heads and Gold Coast area, Australia. Sources spell the name Minjungbal, Minyungbal or Minjunbal. It is treated as a dialect within the wider Yugambeh-Bundjalung group. Search results show a recent scholarly dictionary and grammar and some community pages. They show no Minjungbal-specific school curriculum, primer or audio course. The search results cover English-language and academic materials only; I found no resource written in the language itself beyond word lists.

\section{Primary Educational Resources}
No primers, textbooks or school programmes specific to Minjungbal were found in the search results. The nearest material is the Gurgun Mibinyah dictionary and grammar below, which could serve as a source for teacher-made materials.

\section{Linguistic and Reference Materials}
\begin{itemize}
\item \emph{Gurgun Mibinyah}, a dictionary and grammar of the northern varieties of Yugambeh-Bundjalung (Tweed River to Logan River), compiled and edited by Margaret Sharpe and published in 2020 by Aboriginal Studies Press (AIATSIS). It has English translations, example sentences, a grammar guide and an English word index. The compiler notes that some pronunciations are informed guesses. \url{https://www.nit.com.au/mibiny-languages-immortalised-in-new-dictionary/}
\item Further coverage of the same dictionary: \url{https://www.yarn.com.au/blogs/yarn-in-the-community/the-words-of-the-mibiny-peoples-immortalised-in-new-dictionary}
\item Publisher listing for the dictionary: \url{https://publishers.newsouthbooks.com.au/collections/february-2020-academic-specialist/}
\item Muurrbay Bundjalung online dictionary, entry giving Minyungbal as the language name around Tweed Heads: \url{https://bundjalung.dalang.com.au/language/view_word/3475}
\item Muurrbay Bundjalung online dictionary, entry for Minyungbal dialect(s): \url{https://bundjalung.dalang.com.au/language/view_word/3476}
\end{itemize}

\section{Religious and Community Resources}
\begin{itemize}
\item Tweed Regional Museum, ``Aboriginal cultural heritage'' page, which gives the community's own account of the Tweed dialect (Ngandowal) and the Minyungbal name: \url{https://museum.tweed.nsw.gov.au/explore/aboriginal-cultural-heritage}
\item Coodjingburra community site, ``our story'': \url{https://coodjingburrahome.com/our-story}
\end{itemize}
No scripture translations were found.

\section{Digital and Audio Resources}
The Muurrbay online Bundjalung dictionary (links above) is the only digital resource found. The Gurgun Mibinyah introduction reportedly draws on some later audio recordings, but I found no accessible audio collection.

\section{Curriculum Development Context}
The language is reported as spoken as a first language by the elderly only, and it is not known to be taught in schools. A curriculum would have to be built as a revival project. The base would be the Sharpe dictionary and grammar, with the Tweed and Gold Coast Aboriginal communities in control of content and spelling. Spelling and pronunciation are still debated, so community agreement on orthography should come first. Wider Yugambeh-Bundjalung materials may help but were not verified here.

\end{document}